\chapter{Unitarity Bounds and Short Multiplets}
\label{ch:volume-iii-unitarity-bounds-short-multiplets}

\section{Positivity in Radial Quantization}

Assume the conformal theory is unitary after analytic continuation to
Lorentzian signature and satisfies reflection positivity in Euclidean radial
quantization.  Then the Hilbert space on \(S^{D-1}\) has a positive
semidefinite inner product, and null states are quotiented to obtain the
physical Hilbert space.  The adjoint relation
\[
  \widehat P_\mu^\dagger=\widehat K_\mu
\]
implies that every Gram matrix of descendant states must be positive
semidefinite.

Let \(\ket{\mathcal O_a}\) be a primary of dimension \(\Delta\) and spin
representation \(\rho\).  At level one the descendant Gram matrix is
\begin{equation}
  G_{\mu a,\nu b}^{(1)}
  =
  \bra{\mathcal O_a}\widehat K_\mu\widehat P_\nu\ket{\mathcal O_b}.
  \label{eq:level-one-gram-general}
\end{equation}
Using \(\widehat K_\mu\ket{\mathcal O_b}=0\) and the conformal algebra gives
\[
  G_{\mu a,\nu b}^{(1)}
  =
  -2\ii
  \bra{\mathcal O_a}
  (\delta_{\mu\nu}\widehat D-\widehat J_{\nu\mu})
  \ket{\mathcal O_b}.
\]
This equation is a representation-theoretic constraint.  Positivity of
\(G^{(1)}\) gives lower bounds on \(\Delta\) once the \(SO(D)\) representation
\(\rho\) is specified.

\section{Scalar Primaries}

For a scalar primary \(\phi\), the spin matrices vanish.  The level-one Gram
matrix reduces to
\[
  \bra{\phi}\widehat K_\mu\widehat P_\nu\ket{\phi}
  =
  2\Delta_\phi\delta_{\mu\nu}\braket{\phi}{\phi}.
\]
Thus
\[
  \Delta_\phi\ge0.
\]

The sharper scalar bound comes from level two.  Define
\[
  M^{(2)}_{\alpha\beta,\mu\nu}
  =
  \bra{\phi}
  \widehat K_\alpha\widehat K_\beta
  \widehat P_\mu\widehat P_\nu
  \ket{\phi}.
\]
Repeated use of the conformal commutation relations gives, for a scalar
primary,
\begin{equation}
  M^{(2)}_{\alpha\beta,\mu\nu}
  =
  4\left[
    (\delta_{\alpha\mu}\delta_{\beta\nu}
    +\delta_{\alpha\nu}\delta_{\beta\mu})
    \Delta_\phi(\Delta_\phi+1)
    -
    \delta_{\alpha\beta}\delta_{\mu\nu}\Delta_\phi
  \right]\braket{\phi}{\phi}.
  \label{eq:scalar-level-two-gram}
\end{equation}
Contracting with \(\delta^{\alpha\beta}\delta^{\mu\nu}\) computes the norm of
\(\widehat P^2\ket{\phi}\):
\begin{equation}
  \left\|\widehat P^2\ket{\phi}\right\|^2
  =
  4D\,\Delta_\phi\bigl(2\Delta_\phi+2-D\bigr)
  \braket{\phi}{\phi}.
  \label{eq:scalar-laplacian-descendant-norm}
\end{equation}
For a positive-norm scalar primary, positivity gives the alternatives
\begin{equation}
  \Delta_\phi=0
  \quad\text{with}\quad
  \widehat P_\mu\ket{\phi}=0,
  \qquad
  \text{or}
  \qquad
  \Delta_\phi\ge\frac{D-2}{2}.
  \label{eq:scalar-unitarity-bound}
\end{equation}
The first branch is the identity representation.  The second branch is the
scalar unitarity bound.

When
\[
  \Delta_\phi=\frac{D-2}{2},
\]
the state \(\widehat P^2\ket{\phi}\) is null.  In the quotient Hilbert space
this gives the operator equation
\begin{equation}
  \partial^2\phi(x)=0.
  \label{eq:scalar-shortening-laplace}
\end{equation}
The corresponding conformal multiplet is the free massless scalar multiplet.

\section{Symmetric Traceless Tensor Primaries}

Let
\[
  \mathcal O_{\mu_1\cdots\mu_\ell}(x)
\]
be a primary transforming as a rank-\(\ell\) symmetric traceless tensor of
\(SO(D)\), with \(\ell\ge1\).  Thus the indices are symmetric and
\[
  \delta^{\mu_1\mu_2}
  \mathcal O_{\mu_1\mu_2\cdots\mu_\ell}=0.
\]
Decomposing the level-one descendants into irreducible \(SO(D)\)
representations gives the bound
\begin{equation}
  \Delta_{\mathcal O}\ge \ell+D-2.
  \label{eq:spin-ell-unitarity-bound}
\end{equation}
At equality, the divergence descendant is null:
\begin{equation}
  \partial^{\mu_1}
  \mathcal O_{\mu_1\mu_2\cdots\mu_\ell}(x)=0.
  \label{eq:spin-ell-conservation-shortening}
\end{equation}
For \(\ell=1\) this is current conservation and
\(\Delta_J=D-1\).  For \(\ell=2\) it is stress-tensor conservation and
\(\Delta_T=D\).

\section{Short Multiplets as Operator Equations}

A short conformal multiplet is a primary multiplet for which a descendant has
zero norm and is set to zero in the physical Hilbert space.  In local-operator
language, the null descendant is an operator equation.  The equations
\[
  \partial^2\phi=0,\qquad
  \partial^\mu J_\mu=0,\qquad
  \partial^\mu T_{\mu\nu}=0
\]
are therefore representation-theoretic statements in a unitary conformal
theory.  They also carry dynamical information: a conserved current generates
a global symmetry, the stress tensor generates spacetime symmetries, and a
scalar saturating the scalar bound generates the free scalar multiplet.

The unitarity bounds thus convert positivity of the Hilbert-space inner
product into exact restrictions on the allowed local operator spectrum.
